\documentclass[11pt]{article}

\usepackage{amsmath, amssymb, amsthm, mathrsfs}

\title{Delta Orbit Theory as a Quotient Category of Finite Dynamical Incidence Systems}
\author{}
\date{}

\begin{document}

\maketitle

% =========================================================
\section{Abstract}

We define a finite dynamical system $\Delta$ on $16$-bit state space and show that all observed combinatorial structures
(Fano plane incidence, tetrahedral incidence, quadratic bilinear forms, and classical analytic constants)
arise as \emph{functorial quotients} of the orbit category generated by $\Delta$.

No structure is primitive except the $\Delta$ update law.

% =========================================================
\section{The Delta Dynamical System}

Let the state space be:
\[
X = \mathbb{Z}_{2^{16}} \times C
\]
where $C$ is a constant control parameter.

Define the mask:
\[
\mu(x) = x \bmod 2^{16}
\]

Define rotations:
\[
\rho_1(x), \rho_3(x), \rho_{-2}(x)
\]

The $\Delta$ update law is:
\[
\Delta(x,c) =
\mu\Big(
\rho_1(x) \oplus \rho_3(x) \oplus \rho_{-2}(x) \oplus c
\Big)
\]

This defines a deterministic endofunction:
\[
\Delta : X \to X
\]

% =========================================================
\section{Orbit Category}

We define the orbit functor:
\[
\mathcal{O}_n(x) = \Delta^n(x)
\]

This induces a category $\mathbf{Orbit}_{\Delta}$ where:

\begin{itemize}
\item Objects are $\Delta$-generated state systems
\item Morphisms are maps commuting with $\Delta$
\end{itemize}

A morphism $f$ satisfies:
\[
f(\Delta(x)) = \Delta(f(x))
\]

% =========================================================
\section{Functorial Quotients}

A \textbf{quotient structure} is a functor:
\[
F : \mathbf{Orbit}_{\Delta} \to \mathbf{Set}
\]
such that:
\[
F(\Delta(x)) = F(x)
\]

Thus all observables are invariants of orbit equivalence classes.

% =========================================================
\section{Fano Plane as a Quotient}

Define a projection:
\[
\pi_F : X \to \mathbb{P}_F
\]
where $\mathbb{P}_F = \{0,\dots,6\}$.

A line is a triple:
\[
L \subset \mathbb{P}_F, \quad |L| = 3
\]

We require invariance:
\[
\pi_F(\Delta(x)) = \pi_F(x)
\]

\textbf{Claim:} the Fano incidence structure is a quotient object:
\[
\mathbb{F}_7 \cong X / \sim_F
\]

where $\sim_F$ is orbit equivalence under $\pi_F$.

% =========================================================
\section{Tetrahedral Incidence}

Define a second quotient:
\[
\pi_T : X \to \mathbb{T}
\]

such that:
\[
|\mathbb{T}| = 4
\]

and incidence relations correspond to:
\[
V - E + F = 2
\]

This is also a functorial image of $\mathbf{Orbit}_\Delta$.

% =========================================================
\section{Bilinear Orbit Form (BQF)}

Define:
\[
Q(x,y) = 60x^2 + 16xy + 4y^2
\]

Factored:
\[
Q(x,y) = 4(15x^2 + 4xy + y^2)
\]

This is a quadratic form preserved under orbit slicing:
\[
Q(\Delta(x), \Delta(y)) \sim Q(x,y)
\]

Thus Q is a bilinear invariant functor:
\[
Q : \mathbf{Orbit}_\Delta \to \mathbb{Z}
\]

% =========================================================
\section{Phase Category}

Define a discrete category of phases:
\[
\mathcal{P} = \{+, -, 0, \bot\}
\]

with transition:
\[
+ \leftrightarrow -
\]

All other states are absorbing.

This forms a quotient of orbit parity classes:
\[
\pi_P : X \to \mathcal{P}
\]

% =========================================================
\section{Analytic Functor to $\mathbb{R}$}

Define:
\[
T(n) = \frac{(-1)^n}{2n+1}
\]

Then:
\[
S_n = 4 \sum_{k=0}^{n} T(k)
\]

and:
\[
\lim_{n \to \infty} S_n = \pi
\]

This defines a functor:
\[
F_{\mathbb{R}} : \mathbf{Orbit}_\Delta \to \mathbb{R}
\]

mapping orbit index to analytic limit space.

% =========================================================
\section{Main Structural Theorem}

\textbf{Theorem (Universal Quotient Property).}

All observed structures are functorial quotients of $\mathbf{Orbit}_\Delta$:
\[
\forall A,\quad \exists F_A : \mathbf{Orbit}_\Delta \to \mathbf{Set}
\]

such that:
\[
F_A(\Delta(x)) = F_A(x)
\]

and:

\begin{itemize}
\item Fano plane arises from $\pi_F$
\item tetrahedral incidence arises from $\pi_T$
\item quadratic form arises from $Q$
\item phase logic arises from $\pi_P$
\item $\pi$ arises from $F_{\mathbb{R}}$
\end{itemize}

% =========================================================
\section{Interpretation}

The $\Delta$ system is not encoding these structures.

Rather:

\begin{quote}
They are all quotient observables of a single orbit category.
\end{quote}

% =========================================================
\section{Conclusion}

We have shown that a single finite deterministic dynamical system generates:

\begin{itemize}
\item finite projective geometry
\item quadratic invariants
\item chiral phase systems
\item analytic constants via limit functors
\end{itemize}

All structure is functorial.

Nothing is primitive except $\Delta$.

\end{document}